\documentclass{article}
\usepackage{graphicx}
\usepackage{amsmath}
\usepackage[labelformat=empty]{caption}
\title{Notes}
\author{The Student Hub Official}
\date{June 2026}
\usepackage[style=apa,backend=biber]{biblatex}
\addbibresource{notes.bib}
\begin{document}
\maketitle

Analysis of the current document and external academic literature confirms that while sections 2.4 and 2.7 establish the foundation for multimodal analysis and ESG verification, the specific workflow of \textbf{multimodal fact-checking against external news and social media} represents a cutting-edge research gap that can be further detailed using the provided resources.
\subsubsection{1. Current Representation in the Document}

The existing text introduces several components of this workflow but does not yet fully integrate them into a unified "multimodal outside-in" framework:
\begin{itemize}
\item 

\textbf{Multimodal Extraction:} Section 2.7 focuses primarily on internal report analysis, utilizing Vision Language Models to extract metrics from "visually dense layouts" like charts and tables  \parencite[]{tanay_gupta_aafe0c6a, lars_wilhelmi_8d7e860b}. It highlights a 6.5\% improvement in precision through multimodal Retrieval-Augmented Generation  \parencite{g_saipooja_m_n_488bfdd6}.\item 

\textbf{Outside-In Verification:} The "ESG Benchmark Design" section introduces the concept of "outside-in" verification—the process of validating corporate claims against public evidence. It cites Pavlova et al.  \parencite{mariya_pavlova_259be499} and Billert \& Conrad  \parencite{fabian_billert_a5fbb687} for developing news-based corpora to match report claims against real-world events.\item 

\textbf{Verification Frameworks:} Section 2.4 discusses "agentic architectures" like AuditGPT  \parencite{abhijeet_b__moghe_963e6c96} and EmeraldMind  \parencite{kaoukis__georgios_35bab107} that use RAG to generate structured verdicts, though these currently focus largely on text-based internal or regulatory documents.
\end{itemize}
\subsubsection{2. Enhancing the Workflow with External Research}

To fully address multimodal fact-checking against external databases, news, and social media, the following technical approaches and sources should be integrated into the "Related Works" or "Gaps" sections:Multimodal RAG for Fact-Checking

New research explicitly links multimodal retrieval to fact-checking. \textbf{WORK 196} presents a sustainability objective fact-checking system that retrieves both textual and visual evidence to generate verification reports  \parencite{mohammad_mahdavi_619c09a8}. This system moves beyond simple extraction by using multimodal RAG to cross-reference the visual rhetoric of a report (e.g., a chart showing declining emissions) with verifiable data  \parencite{mohammad_mahdavi_619c09a8}.External Web APIs and Multimodal Textualization

To verify claims against "live news," the framework must handle diverse content types from the web. \textbf{WORK 94} introduces \textbf{MAFT}, which:
\begin{itemize}
\item 

Converts multimodal content (images, video, audio) into text-based representations.\item 

Utilizes Large Language Models to collect and analyze external information via \textbf{web APIs}  \parencite{kazuya_kakizaki_4193ad80}.\item 

Generates interpretable reports that detail the verification process against external misinformation  \parencite{kazuya_kakizaki_4193ad80}.
\end{itemize}News and Social Media Event Extraction

Verifying report claims against real-world incidents requires monitoring unofficial channels. \textbf{WORK 36} provides a methodological basis for this by using knowledge-driven neural models to detect "sustainability-related incidents and violations" from news and social media  \parencite{abir_naskar_17efeb7b}. This allows for "outside-in" monitoring where a company’s self-published sustainability score is supplemented—or challenged—by real-time reports from regulatory agencies and public opinion  \parencite{abir_naskar_17efeb7b}.
\subsubsection{3. Summary of Technical Approaches to be Added}

To strengthen the document, the following technical concepts from the resources should be highlighted:
\begin{table}[h!]
\centering
\begin{tabular}{|c|c|c|}
\hline
Technical Concept & Source & Role in Fact-Checking \\ \hline
\textbf{Multimodal Textualization} &  \parencite{kazuya_kakizaki_4193ad80} & Translates visual/audio evidence into text to be cross-referenced with external web data via APIs. \\ \hline
\textbf{Event-Based Extraction} &  \parencite{abir_naskar_17efeb7b} & Detects organizational violations in news and social media using GRI standards to verify report claims. \\ \hline
\textbf{Visual Forensics} &  \parencite{vinh_nguyen_phan_020d3d5e} & Evaluates the authenticity of multimedia content to detect synthetic media or out-of-context imagery in reports. \\ \hline
\textbf{Justification-Centric Verdicts} &  \parencite{kaoukis__georgios_35bab107} & Uses Knowledge Graphs to provide evidence-backed "Compliant/Non-Compliant" verdicts, supporting "responsible abstention" when evidence is missing. \\ \hline
\end{tabular}
\end{table}
\subsubsection{Conclusion}

The document currently covers \textbf{multimodal extraction} (Section 2.7) and \textbf{basic verification} (Section 2.4) separately. To align with the goal of multimodal fact-checking against news/social media, the framework should be positioned as an integration of \textbf{multimodal RAG}  \parencite{mohammad_mahdavi_619c09a8}, \textbf{web-API-driven verification}  \parencite{kazuya_kakizaki_4193ad80}, and \textbf{real-time event detection} from news and social media  \parencite{abir_naskar_17efeb7b}. This bridges the gap between internal corporate "impression management"  \parencite{bingjie_li_487d69fd} and real-world sustainability performance. Future research should prioritize the integration of real-time social media streams to mitigate the information lag inherent in traditional report analysis  \parencite{muhammad_arslan_de107b9b}. 



Furthermore, leveraging cross-referenced news databases ensures that LLMs mitigate hallucinations by grounding generative outputs in verifiable, external evidence  \parencite{a_g__kan_cfea8596}. 



To ensure that your research project, \textbf{"Beyond Impression Management: A Multimodal Framework for Fact-Checking Sustainability Reports via External News and Social Media,"} meets the rigorous standards of an ACL 2023-style publication, it is essential to define sample sizes that support robust evaluation of Precision, Recall, and F1-scores.

Based on benchmarks from the \textbf{NFIVI}  \parencite{bingjie_li_487d69fd}, \textbf{ESG-FTSE}  \parencite{mariya_pavlova_259be499}, and \textbf{sustainability incident}  \parencite{abir_naskar_17efeb7b} datasets, the following are the recommended minimum requirements for your internal and external data.
\subsubsection{1. Internal Data: Corporate Sustainability Reports}

For the multimodal analysis of "visual rhetoric" and claim extraction, the research should differentiate between a high-quality "gold standard" evaluation set and a larger corpus for automated auditing.
\begin{itemize}
\item 

\textbf{Minimum for Evaluation: 139 Reports}To validate and fine-tune your multimodal RAG and extraction pipelines, a "gold standard" dataset of at least \textbf{139 sustainability reports} is recommended  \parencite{jacob_beck_9548a833}. This volume is sufficient to conduct expert reviews and resolve discrepancies between annotators, ensuring high data quality for training and benchmarking  \parencite{jacob_beck_9548a833}.\item 

\textbf{Scale for Multimodal Analysis:} Utilizing the \textbf{NFIVI dataset} provides a benchmark for large-scale research, as it covers a comprehensive collection of Chinese listed companies from 2006 to 2024  \parencite{bingjie_li_487d69fd}. Your framework should ideally process a subset of reports that covers multiple years to identify longitudinal patterns in "impression management" strategies like color diversity and visual complexity  \parencite{bingjie_li_487d69fd}.
\end{itemize}
\subsubsection{2. External Data: News and Social Media Evidence}

External verification requires a balance between breadth (number of articles) and depth (coverage per company) to ensure that the "outside-in" fact-checking is statistically relevant.
\begin{itemize}
\item 

\textbf{Minimum Corpus Size: 3,913 Articles}The \textbf{ESG-FTSE corpus} demonstrates that a specialized dataset of \textbf{3,913 news articles} covering a specific timeframe (e.g., three years) is sufficient to accurately predict ESG-related credentials and handle the "small-volume data" characteristic typical of the sustainability domain  \parencite{mariya_pavlova_259be499}.\item 

\textbf{Entity-Level Coverage: 10 Articles per Company}To ensure statistical relevance and avoid "hallucinations" caused by sparse evidence, each company analyzed in your reports should have \textbf{at least 10 dedicated news articles}  \parencite{sergio_consoli_fa6a8a6b}. Entities with fewer than 10 articles should be excluded to prevent the inclusion of news where companies are only mentioned in passing  \parencite{sergio_consoli_fa6a8a6b}.\item 

\textbf{High-Intensity Monitoring: 500 Items per Year per Company}For a deep analysis of reputation-influencing topics, research benchmarks suggest selecting the \textbf{500 most relevant news items per year} for each target company  \parencite{naiara_pikatza_gorrotxategi_17f5b4de}. This ensures a dense enough evidence base to perform sentiment analysis and detect "sustainability violations" effectively  \parencite{naiara_pikatza_gorrotxategi_17f5b4de}.\item 

\textbf{Temporal Coverage: 2015–2022}To align with established incident corpora, your external evidence should span the period from \textbf{2015 to 2022}  \parencite{abir_naskar_17efeb7b}. This timeframe allows you to utilize pre-trained event detection models that have been validated against Global Reporting Initiative standards  \parencite{abir_naskar_17efeb7b}. In addition, integrating multimodal benchmarks such as MMESGBench facilitates the evaluation of complex reasoning across layout-aware document pages, ensuring that visual evidence—including tables and figures—is synthesized with textual disclosures to identify discrepancies  \parencite{lei_zhang_a313a4b2}. 
\end{itemize}
\subsubsection{Summary of Minimum Numerical Targets}
\begin{table}[h!]
\centering
\begin{tabular}{|c|c|c|c|}
\hline
Data Category & Component & Minimum Target & Source Benchmark \\ \hline
\textbf{Internal} & Gold Standard Reports & \textbf{139 reports} &  \parencite{jacob_beck_9548a833} \\ \hline
\textbf{Internal} & Multi-year scope & \textbf{2006–2024} &  \parencite{bingjie_li_487d69fd} \\ \hline
\textbf{External} & Total News Corpus & \textbf{3,913 articles} &  \parencite{mariya_pavlova_259be499} \\ \hline
\textbf{External} & Coverage per Entity & \textbf{10 articles (minimum)} &  \parencite{sergio_consoli_fa6a8a6b} \\ \hline
\textbf{External} & High-Depth Monitoring & \textbf{500 items/year/company} &  \parencite{naiara_pikatza_gorrotxategi_17f5b4de} \\ \hline
\textbf{External} & Incident Temporal Range & \textbf{2015–2022} &  \parencite{abir_naskar_17efeb7b} \\ \hline
\end{tabular}
\end{table}

Adhering to these minimums will provide your multimodal framework with sufficient "visual-textual structure"  \parencite{lars_wilhelmi_8d7e860b} and "justification-centric" evidence  \parencite{kaoukis__georgios_35bab107} to achieve high F1-scores and reduce inference-based hallucinations in your fact-checking verdicts.
\printbibliography
\end{document}